\newpage\section{What follows for coefficients and cohomology}

Write \(k=\overline{\mathbb F}_p\), \(N=2m+1\), and \(G=\mathrm{SL}_N(k)\).
For degree one, both papers assume Lusztig's character formula in the
given rank and \(p\ge5(N-1)\) for finite group comparison. These assumptions
hold for every sufficiently large \(p\) at fixed \(m\).

\begingroup\small\renewcommand{\arraystretch}{1.17}
\begin{tabularx}{\linewidth}{@{}>{\raggedright\arraybackslash}p{.215\linewidth}Y@{}}
\toprule
Setting and source & Consequence \\ \midrule
A, Theorem 4L.1 and Corollary 6C.2 &
\(\sum_{j=1}^m\mu_{j;m}=m(m+1)/2\).
Each coefficient equals the rational \(H^1\)-dimension of
\(L(\lambda_{j;m}(p))\) and the finite \(H^1\)-dimension for its restriction
to \(\mathrm{SL}_N(\mathbb F_p)\) and its descendant \(M_{j;m}(p)\)
on \(\mathrm{PSL}_N(\mathbb F_p)\), over \(k\).
The row sum gives \(\max_j\mu_{j;m}\ge\lceil(m+1)/2\rceil\).\\
B, Theorems A and B &
\(\mu_{m;m}=m\), and
\(\dim_{\mathbb F_p}H^1(\mathrm{PSL}_N(\mathbb F_p),V_{m,p})=m\).
All \(V_{m,p}\) are faithful and absolutely irreducible.\\
The two KL theorems together &
\(\sum_{j<m}\mu_{j;m}=m(m-1)/2\) and \(\max_j\mu_{j;m}\ge m\).\\
Partial-row calculation, \S\ref{sec:partialrows} &
$\mu_{j;m}=j$ throughout, with maximum $m$ at $j=m$. For each sufficiently large $p$, our $m$ modules realize every first cohomology dimension $1,\ldots,m$ for this one finite group.\\
B, Theorem B, (2.17), Corollary 2.20 &
For \(0\le s<m\), the normalized polynomial for \(y_s\) starts with
\(mv^{s+1}\). For large \(p\),
\(\dim_k H^{s+1}(G,L(y_s\cdot_p0))=m\), and its rational
cohomology vanishes in all lower degrees. The identified diagonal is \(s=0\).\\
B, Theorem C, Lemma 2.18 &
Fix \(1\le n\le m\). For sufficiently large \(p\), then sufficiently
large \(a\), a faithful absolutely irreducible \(k\)-module for
\(\mathrm{PSL}_N(\mathbb F_{p^a})\) has \(\dim_kH^n\ge m\).
The proof uses injective restriction.\\
B, Theorem D, (2.18) &
For a primitive complex root of odd order \(\ell>N\),
\(\dim_{\mathbb C}H^{s+1}(U_\zeta,L_\zeta(y_s\cdot_\ell0))=m\),
with vanishing in lower degrees. This is Ext in finite-dimensional
integrable type-1 modules of Lusztig's divided-power quantum group.\\
A with B's (2.18) &
The same standard quantum KL formula transports our row sum to quantum
\(H^1\) for the \(m\) modules \(L_\zeta(u_{j;m}^{-1}\cdot_\ell0)\).\\
\bottomrule
\end{tabularx}\endgroup

\paragraph{The coefficient fields.}
On the diagonal the highest weights agree. The split group gives an
\(\mathbb F_p\)-form of the simple highest-weight module. Its root-lattice
weight makes the finite center act trivially. That center has
prime-to-\(p\) order, so inflation across it is an isomorphism in every
degree. Thus, as modules for the \emph{same} finite group,
\[
 M_{m;m}(p)\cong V_{m,p}\otimes_{\mathbb F_p}k.
\]
Extending the finite group bar complex shows that the scalar extension
preserves the cohomology dimension. The same descent applies to our
other restricted highest weights. See A Lemmas 6A.2, 6B.1 and 6C.1;
B Proposition 2.15 and the proof of Theorem A.

\paragraph{Finite-group comparison in higher degrees.}
For the weights in the parabolic paper, the standard large-characteristic comparison gives exact
finite-group cohomology dimensions.
Fix \(1\le n\le m\), set \(\lambda=y_{n-1}\cdot_p0\), and assume
the character-formula hypotheses in that paper together with
\(p\ge(4n+1)(N-1)\).
B Lemma~2.17 makes \(\lambda\) nonzero, \(p\)-restricted and a root-lattice
weight in the affine dot orbit of zero. It is therefore \(p^a\)-restricted
for every \(a\ge1\).
Parshall--Scott--Stewart, \emph{Shifted generic cohomology},
Theorem~5.12 and its proof on pp.~1780--1781, give
\[
 H^d(\mathrm{SL}_N(\mathbb F_{p^a}),L(\lambda))
 \cong H^d(G,L(\lambda))
 \qquad(1\le d\le n,\ a\ge1).
\]
Here one takes \(\nu=0\) in the Bendel--Nakano--Pillen
Theorem~7.5 comparison displayed in that proof.
Combining this with Theorem~B and (2.17) of the parabolic paper, and inflating across the
prime-to-\(p\) center, gives a faithful absolutely irreducible
\(k\mathrm{PSL}_N(\mathbb F_{p^a})\)-module \(V_a\) with
\[
 \dim_k H^n(\mathrm{PSL}_N(\mathbb F_{p^a}),V_a)=m,\qquad
 H^d(\mathrm{PSL}_N(\mathbb F_{p^a}),V_a)=0\quad(0\le d<n).
\]
Degree zero vanishes by nontrivial irreducibility.
In particular, equality and lower-degree vanishing hold for \(a=1\).
This follows from the rational cohomology calculation in the parabolic paper and the cited
comparison theorem.
The required lower bound on the characteristic depends on the rank and
the chosen degree.
\href{https://doi.org/10.1112/S0010437X13007331}{doi:10.1112/S0010437X13007331}.
